% Start preamble
\documentclass[12pt,a4paper]{article}
\usepackage{geometry}
 \geometry{
 a4paper,
 total={170mm,257mm},
 left=20mm,
 top=20mm,
 }
\usepackage{currfile}
\usepackage[utf8]{inputenc}
\usepackage[T1]{fontenc}
\usepackage[pdftex]{graphicx}
\graphicspath{{./}}
\usepackage{enumitem}
\usepackage{pdfpages}
\usepackage{hyperref}
\usepackage{tikz}
\usepackage{attachfile}
\usepackage{epstopdf}
\usepackage{array}
\usepackage{multirow}
\usepackage{multicol}
\usepackage{float}
%\usepackage[table]{xcolor,colorbl}
\setlength{\textwidth}{16cm}
\setlength{\oddsidemargin}{-0.5cm}
\setlength{\evensidemargin}{-0.5cm}
%\setlenght{\headsep}{0cm}
\setlength\parindent{0pt}
%\setlength{\extrarowheight}{3pt}
\usepackage{listings}
%\usepackage{xcolor}
\input{arduinoLanguage.tex}
%%%%%% Counting oppgaves %%%%%%
 \newcount\questnum \questnum=0
 \def\oppgave{
            \advance\questnum by 1
            \ifnum \questnum > 0
                 \hrule
                 \vskip 5pt
                 \leftline{Oppgave \the\questnum}
                 \vskip 5pt \fi
		}
% End preamble
\begin{document}
\title{Oppgaver til sikkerhetsrelaterte deler av styresystemer}
\author{Faglærer: Fred-Olav Mosdal 90507684\\}
\maketitle
\oppgave{} 
Hvilke to standarder finnes for funksjonellsikkeret i maskiner og hvilke deler av styresystemet på en maskin kan de brukes til?
% 
\vskip 2.5pt 
\begin{tikzpicture}
	\draw[step=0.5cm,gray!20,very thin]  grid (17,8) ;
\end{tikzpicture}
\vskip 5pt 
\oppgave{} 
Hva står PL og SIL for og hva forteller disse?
%
\vskip 2.5pt 
\begin{tikzpicture}
	\draw[step=0.5cm,gray!20,very thin]  grid (17,4) ;
\end{tikzpicture}
\vskip 5pt 
\oppgave{} 
I teksten beskrives det at en standard er harmonisert under maskindirektivet, hva betyr dette? Søk på internett.

%
\vskip 2.5pt 
\begin{tikzpicture}
	\draw[step=0.5cm,gray!20,very thin]  grid (17,4) ;
\end{tikzpicture}
\vskip 5pt 
\oppgave{} 
Hvilken standard brukes av flest maskinprodusenter?
%
\vskip 2.5pt 
\begin{tikzpicture}
	\draw[step=0.5cm,gray!20,very thin]  grid (17,4) ;
\end{tikzpicture}
\vskip 5pt 
\oppgave{} Er en sikkerhetsfunksjon i styresystemet første valget ved konstruksjon av maskiner?
%Nei første prioritet ved design av an maskin er å konstruere den slik at det ikke er behov for sikkerhetsfunksjoner. Om dette ikke er mulig bygges vern mot farene. Sikkerhetsfunksoner i styressystemet er ofte en del av dette vernet. F.eks. en bevegelig dør. 

\vskip 2.5pt 
\begin{tikzpicture}
	\draw[step=0.5cm,gray!20,very thin]  grid (17,4) ;
\end{tikzpicture}
\vskip 5pt 
\oppgave{} Hva er en sikkerhetsfunksjon?
\vskip 2.5pt 
\begin{tikzpicture}
	\draw[step=0.5cm,gray!20,very thin]  grid (17,4) ;
\end{tikzpicture}
\vskip 5pt 
%En sikkerhetsfunksjon i styresystemet er en funksjon som skal bidra til å redusere risikoen ved en identifisert fare med maskinen

\oppgave{} Hva står PL for?
\vskip 2.5pt 
\begin{tikzpicture}
	\draw[step=0.5cm,gray!20,very thin]  grid (17,4) ;
\end{tikzpicture}
\vskip 5pt 

%PL står for Preformance Level dette forteller noen om sansynligheten for at en sikkerhetsfunksjon virker når det er behov for en. 

	
	

\vskip 10pt 
Oppgaver fra Kap 5. 



	\oppgave{} Gjelder EC machinery Directive i norge? Eventuelt hvilken lov er det implementert i?
\vskip 2.5pt 
\begin{tikzpicture}
	\draw[step=0.5cm,gray!20,very thin]  grid (17,4) ;
\end{tikzpicture}
\vskip 5pt 
%Den gjelder i norge som en del av EØS avtalen og er implementert i forskrit om maskiner. 


	\oppgave{} Hvilken sikkerhetsstandard er viktig grunnlag for maskin direktivet?
\vskip 2.5pt 
\begin{tikzpicture}
	\draw[step=0.5cm,gray!20,very thin]  grid (17,4) ;
\end{tikzpicture}
\vskip 5pt 
%EN ISO 12100. Dette er en veldig generell standard som brukes for utgangpunkt for alle. Det finnes også B type standarder som omhandler spesiele deler av en maskin f.eks. EN 60204-1 som omhandler elektrisk utstyr i maskiner eller EN ISO 13849-1 som gjelder sikkerhetsrelaterte deler av styresystemer. Den siste standarden er C type standarder som gjelder for spesielle maskintyper f. eks. roboter. 

	\oppgave{} Hva er en Type C standard under maskindirektivet?
\vskip 2.5pt 
\begin{tikzpicture}
	\draw[step=0.5cm,gray!20,very thin]  grid (17,4) ;
\end{tikzpicture}
\vskip 5pt 
%Det er en type standard som gjeler en maskin type. Når du skal konstruer en maskin må en alltid først sjekke om det finnes en C-type standard å følge. 

	\oppgave{} Beskriv prosedyren for risikoreduksjon av en maskin. 
\vskip 2.5pt 
\begin{tikzpicture}
	\draw[step=0.5cm,gray!20,very thin]  grid (17,4) ;
\end{tikzpicture}
\vskip 5pt 
%		○ Bestemm grensene til maskinen
%		○ Identifiser farene
%		○ Estimer risikoen
%		○ Eveluer risikoen
%		○ Er risikoen lav nok? Ja/Nei (Ja slutt)
%		
		

	\oppgave{} Hvordan finner en frem til nødvendig sikkerhetsfunksjoner
\vskip 2.5pt 
\begin{tikzpicture}
	\draw[step=0.5cm,gray!20,very thin]  grid (17,4) ;
\end{tikzpicture}
\vskip 5pt 
%For hver fare som er identifisert kommer en opp med en sikkerhetsfunksjon f.eksk. Vern og så vurderer en på nytt om risikoen er tilstrekkelig redusert. 

	\oppgave{} Hvilke sikkerhetsfunksjoner er definert i EN ISO 13849-1
\vskip 2.5pt 
\begin{tikzpicture}
	\draw[step=0.5cm,gray!20,very thin]  grid (17,4) ;
\end{tikzpicture}
\vskip 5pt 
%		○ Sikkerhetsrelatert stoppfunksjon fra et bevegelig vern
%		○ Manuell resett
%		○ Start/restart
%		○ Manuell styring i faresone
%		○ Hold ut run
%		○ Enable (f.eks. knappen på sigen av robotkontrolleren)
%		○ Beskyttelse mot uventetoppstart ved operatør arbeid i faresone
%		○ Funsjon for å frigjøre personer som sitter fast. 
%		○ Isolasjon av frakobling av energi (elektrisk, mekanisk, hydraulisk, pneumatisk)
%		○ Styremodus. 
%		○ Utførelse av nødstopp. 

	\oppgave{} Hvordan bestemmer PLr for en sikkerhetsfunksjon. 
\vskip 2.5pt 
\begin{tikzpicture}
	\draw[step=0.5cm,gray!20,very thin]  grid (17,4) ;
\end{tikzpicture}
\vskip 5pt 

\end{document}
